% ==============================================================================
% METODOLOGÍA
% ==============================================================================
% Instrucciones de la plantilla:
% - Presentar de manera clara y estructurada el proceso seguido durante la 
%   práctica de laboratorio.
% - Comenzar con la configuración del entorno de trabajo (herramientas, lenguajes,
%   bibliotecas: OpenGL, GLU, Assimp, GLEW, GLFW, GLM, etc.).
% - Detallar los algoritmos gráficos implementados explicando su fundamento teórico
%   y su adaptación práctica (dibujo de formas, transformaciones geométricas, etc.).
% - Incluir imágenes con títulos y descripciones de los pasos ejecutados.
% - Suficientemente detallado para que cualquier lector pueda replicar el programa.
% - REGLA: "Una metodología por actividad".
% ==============================================================================

\section{Metodología}

\subsection{Configuración del entorno de trabajo}
Para el desarrollo de la práctica se utilizó un entorno de desarrollo estructurado en C++ configurado con soporte para OpenGL moderno. En la Tabla \ref{tab:requerimientos} se detallan las herramientas, bibliotecas y componentes del sistema empleados.

\begin{table}[H]
\centering
\renewcommand{\arraystretch}{1.25}
\begin{tabular}{|p{4.2cm}|p{10.5cm}|}
\hline
\textbf{Elemento} & \textbf{Especificación / Versión empleada} \\ \hline
Sistema operativo & Linux / Windows 11 (64 bits) \\ \hline
Entorno de desarrollo & Visual Studio Code / Microsoft Visual Studio \\ \hline
Compilador y estándar & GCC / MSVC (soporte para \texttt{C++20}) \\ \hline
API gráfica & OpenGL 3.3 (Perfil Core) \\ \hline
Bibliotecas principales & GLFW (manejo de ventanas y eventos), GLAD (cargador de funciones), GLM (matemáticas vectoriales y matriciales), Assimp (carga de modelos 3D) \\ \hline
Hardware de prueba & Tarjeta gráfica compatible con OpenGL 3.3 Core \\ \hline
\end{tabular}
\caption{Requerimientos de hardware, software y bibliotecas utilizadas.}
\label{tab:requerimientos}
\end{table}

\subsection{Actividad 1: Modelos de iluminación integrados en una Escena}

En esta actividad hicimos todo el trabajo dentro de Microsoft Visual Studio usando C++. La estructura del programa se divide en dos partes: la función \texttt{Start()}, donde cargamos todo lo necesario antes de iniciar, y la función \texttt{Update()}, que corre en un ciclo para dibujar cada fotograma de la escena.

Para compilar el proyecto nos apoyamos en varias librerías. Usamos GLFW y GLAD para abrir la ventana y manejar los controles, GLM para las operaciones matemáticas y Assimp para cargar los modelos 3D (\texttt{House02.fbx} para la casa y \texttt{lightDummy.fbx} para representar el foquito de luz). También agregamos la librería \texttt{irrKlang} para poner música de fondo (\texttt{EternalGarden.mp3}) y la clase \texttt{CubeMap} para ponerle un mapa de entorno al fondo.

Durante la práctica fuimos comentando y descomentando varias líneas de código para ver cómo cambiaba la escena. Por ejemplo, probamos cambiar el modo de dibujo a malla de alambre con las teclas, activamos las transparencias con \texttt{GL\_BLEND} y fuimos ajustando las variables de la luz y del material para ver en tiempo real cómo afectaban el aspecto visual de la casa.

\subsubsection{Algoritmo y fundamento matemático}

El modelo de iluminación de Phong se encarga de calcular cómo se ve la luz al chocar contra la superficie de un objeto, combinando tres partes: ambiental (la luz básica del entorno), difusa (la luz directa que le da a las caras) y especular (el brillo que refleja). La fórmula matemática para calcular el color final $L_o$ suma estas tres componentes:

\begin{equation}
    L_o = \left[ k_a + k_d (\mathbf{n} \cdot \mathbf{l}) + k_s (\mathbf{v} \cdot \mathbf{r})^q \right] \frac{L_i}{r^2}
    \label{eq:phong}
\end{equation}

Donde $k_a$, $k_d$ y $k_s$ son los colores/propiedades del material, $\mathbf{n}$ es la normal de la cara, $\mathbf{l}$ apunta hacia la luz, $\mathbf{v}$ va hacia la cámara y $q$ determina qué tan concentrado se ve el brillo especular.

Además, para acomodar los objetos en el espacio 3D y proyectarlos en la pantalla, multiplicamos la posición de cada punto $P$ por las matrices de Modelo ($M$), Vista ($V$) y Proyección ($P$):

\begin{equation}
    P' = M_{\text{proyección}} \cdot M_{\text{vista}} \cdot M_{\text{modelo}} \cdot P
    \label{eq:transformacion}
\end{equation}

\subsubsection{Implementación del código}

A continuación mostramos un fragmento del código principal dentro de la función \texttt{Update()}, donde activamos el shader de Phong, le pasamos la posición de la cámara y de la luz, ajustamos los colores del material y mandamos a dibujar la casa:

\begin{lstlisting}[language=C++, caption={Configuración de la luz, el material y dibujado del modelo en C++.}, label={lst:actividad1}]
// Activamos el shader de Phong y habilitamos la mezcla de transparencias
phongShader->use();
glEnable(GL_BLEND);
glBlendFunc(GL_SRC_ALPHA, GL_ONE_MINUS_SRC_ALPHA);

// Pasamos las matrices de Proyeccion y Vista
phongShader->setMat4("projection", projection);
phongShader->setMat4("view", view);

// Acomodamos el modelo en el espacio (posicion, rotacion y escala)
glm::mat4 model = glm::mat4(1.0f);
model = glm::translate(model, glm::vec3(0.0f, 0.0f, 0.0f));
model = glm::rotate(model, glm::radians(-90.0f), glm::vec3(1.0f, 0.0f, 0.0f));
model = glm::scale(model, glm::vec3(1.0f, 1.0f, 1.0f));
phongShader->setMat4("model", model);

// Le enviamos al shader las propiedades de la luz y la camara
phongShader->setVec4("LightColor", light.Color);
phongShader->setVec4("LightPower", light.Power);
phongShader->setInt("alphaIndex", light.alphaIndex);
phongShader->setFloat("distance", light.distance);
phongShader->setVec3("lightPosition", light.Position);
phongShader->setVec3("eye", camera.Position);

// Propiedades del material (ambiental, difusa, especular y transparencia)
phongShader->setVec4("MaterialAmbientColor", material.ambient);
phongShader->setVec4("MaterialDiffuseColor", material.diffuse);
phongShader->setVec4("MaterialSpecularColor", material.specular);
phongShader->setFloat("transparency", material.transparency);

// Finalmente dibujamos el modelo 3D
basemodel->Draw(*phongShader);
\end{lstlisting}

\subsection{Actividad 2: [Nombre de la Actividad 2]}
% Describir aquí la metodología correspondiente a la Actividad 2.
% Recordar: "Una metodología por actividad".
\textit{Describir aquí el proceso metodológico de la segunda actividad...}
